% ===== tr-02 块1：第 3 章 pp.8--10（原文页 8--10）=====
% [原文 p.8]
因此 $\iota(\partial/\partial y_j)\omega$ 定义了 $H^1(X_b,\R)$ 中的上同调类 $\gamma_j^*$。类似地，令
\begin{equation*}
g_{ij}=-\int_{\Gamma_i\cap X_b}\iota(\partial/\partial y_j)\operatorname{Im}\Omega;
\end{equation*}
则 $-\iota(\partial/\partial y_j)\operatorname{Im}\Omega$ 定义了 $H^{n-1}(X_b,\R)$ 中的上同调类 $\sum_ig_{ij}\Gamma_i^*$，并且 $\lambda_i=\sum_jg_{ij}dy_j$。于是
\begin{align*}
g(\partial/\partial y_j,\partial/\partial y_k)&=-\int_{X_b}\iota(\partial/\partial y_j)\omega\wedge\iota(\partial/\partial y_k)\operatorname{Im}\Omega\\
&=g_{jk}.
\end{align*}

% [Tutor Note: 这里 $g_{ij}$ 的两次出现要分清：定义式 $g_{ij}=-\int_{\Gamma_i\cap X_b}\iota(\partial/\partial y_j)\operatorname{Im}\Omega$ 是「用 $\Gamma_i$ 这族 $(n-1)$-循环把 $-\iota(\partial/\partial y_j)\operatorname{Im}\Omega$ 的周期展开」的系数；而 $\gamma_j^*,\Gamma_i^*$ 分别是 $H_1$ 与 $H^{n-1}$ 的对偶基，二者按 Poincaré 对偶 $(\gamma_i\cap X_b)\cdot(\Gamma_j\cap X_b)=\delta_{ij}$ 配对。]

另一方面，设 $\check y_1,\dots,\check y_n$ 是使 $d\check y_i=\lambda_i$ 的坐标。则
\begin{equation*}
\partial\check y_i/\partial y_j=g_{ij}=g_{ji}=\partial\check y_j/\partial y_i,
\end{equation*}
于是 $\sum\check y_idy_i$ 是闭 1-形式。从而局部地存在函数 $K$，使得 $\partial K/\partial y_i=\check y_i$ 且 $\partial^2K/\partial y_i\partial y_j=g(\partial/\partial y_i,\partial/\partial y_j)$。简单计算还可确认 $\partial\check K/\partial\check y_i=y_i$。另一方面，
\begin{align*}
g(\partial/\partial\check y_i,\partial/\partial\check y_j)&=g\Big(\sum_k\frac{\partial y_k}{\partial\check y_i}\frac{\partial}{\partial y_k},\sum_l\frac{\partial y_l}{\partial\check y_j}\frac{\partial}{\partial y_l}\Big)\\
&=\sum_{k,l}\frac{\partial y_k}{\partial\check y_i}\frac{\partial y_l}{\partial\check y_j}g(\partial/\partial y_k,\partial/\partial y_l)\\
&=\sum_{k,l}\frac{\partial y_k}{\partial\check y_i}\frac{\partial y_l}{\partial\check y_j}\frac{\partial\check y_k}{\partial y_l}\\
&=\frac{\partial y_j}{\partial\check y_i}=\frac{\partial^2\check K}{\partial\check y_i\partial\check y_j}.
\end{align*}
\proofend
\end{proof}

% [Tutor Note: 末行结论 $\partial^2\check K/\partial\check y_i\partial\check y_j$ 的分母带 check，与 p.7 上 Proposition 2.2 末式（分母不带 check）形成对照；后者已在 p.7 的批注中记为疑似笔误，此处可作旁证。]

于是我们引入带多值凸函数的仿射流形的 Legendre 变换这一概念。

\begin{definition}{3.2.3}
设 $B$ 是仿射流形。$B$ 上的一个\emph{多值函数} $K$ 是开覆盖 $\{(U_i,K_i)\}$ 上的一族函数，使得在 $U_i\cap U_j$ 上 $K_i-K_j$ 是仿射线性的。若 Hessian $(\partial^2K_i/\partial y_j\partial y_k)$ 对所有的 $i$、在任意（等价地，在所有）仿射坐标系 $y_1,\dots,y_n$ 下都是正定的，则称 $K$ 是\emph{凸的}。

给定仿射流形与凸多值函数构成的一对 $(B,K)$，$(B,K)$ 的\emph{Legendre 变换}是一个对 $(\check B,\check K)$，其中 $\check B$ 是 $B$ 的底流形上的一个仿射结构，其坐标局部地由 $\check y_i=\partial K/\partial y_i$ 给出；而 $\check K$ 由
\begin{equation*}
\check K_i(\check y_1,\dots,\check y_n)=\sum\check y_jy_j-K_i(y_1,\dots,y_n)
\end{equation*}
定义。
\end{definition}

% [Tutor Note: 多值函数之所以「多值」，在于 $K_i-K_j$ 只差仿射线性函数，于是 $\check y_i=\partial K/\partial y_i$ 在坐标卡之间只差一个常向量，故仍给出良定义的仿射坐标（即 $\check B$ 的仿射结构）。]

\begin{exercise}{3.2.4}
验证 $\check K$ 也是凸的，并且 $(\check B,\check K)$ 的 Legendre 变换是 $(B,K)$。
\end{exercise}

% [原文 p.9]
\section{半平坦镜像对称}

现在我们暂时忘掉特殊拉格朗日纤维化。转而看看 $B$ 上找到的结构如何给出镜像对称的一个玩具版本。

\begin{definition}{3.3.1}
设 $B$ 是热带仿射流形。
\begin{enumerate}
\item 定义 $\Lambda\subseteq T_B$ 为由 $\partial/\partial y_1,\dots,\partial/\partial y_n$ 局部生成的格局部系统，其中 $y_1,\dots,y_n$ 是局部仿射坐标。这是良定义的，因为转移映射属于 $\R^n\rtimes GL_n(\Z)$。令
\begin{equation*}
X(B):=T_B/\Lambda;
\end{equation*}
这是 $B$ 上的一个环面丛。此外，$X(B)$ 带有如下局部定义的复结构。设 $U\subseteq B$ 是带有仿射坐标 $y_1,\dots,y_n$ 的开集，于是 $T_U$ 有坐标函数 $y_1,\dots,y_n$、$x_1=dy_1,\dots,x_n=dy_n$。则
\begin{equation*}
q_j=e^{2\pi i(x_j+iy_j)}
\end{equation*}
给出 $T_U/\Lambda|_U$ 上的一组全纯坐标，且所诱导的复结构与仿射坐标的选取无关。

稍后我们需要它的一个变体：对 $\epsilon>0$，令
\begin{equation*}
X_\epsilon(B):=T_B/\epsilon\Lambda;
\end{equation*}
它带有由
\begin{equation*}
q_j=e^{2\pi i(x_j+iy_j)/\epsilon}
\end{equation*}
给出的坐标所确定的复结构。（我们稍后会看到，极限 $\epsilon\to0$ 对应于大复结构极限。）

\item 定义 $\check\Lambda\subseteq T_B^*$ 为由 $dy_1,\dots,dy_n$ 局部生成的格局部系统，其中 $y_1,\dots,y_n$ 是局部仿射坐标。令
\begin{equation*}
\check X(B):=T_B^*/\check\Lambda.
\end{equation*}
当然，$T_B^*$ 带有典范辛结构，并且这一辛结构下降到 $\check X(B)$ 上。\proofend
\end{enumerate}
\end{definition}

% [Tutor Note: 原文在 Definition 3.1 的 (2) 之后、段末右端印了一个 □（Halmos 记号，通常用于证明结尾）。这在定义中并不常见，但确系原文如此（p.9，y≈465pt，700 dpi 裁图复核）；译文照原文补出该记号。]

% [Tutor Note: $X_\epsilon(B):=T_B/\epsilon\Lambda$ 只是把格按 $\epsilon$ 缩放，复坐标变为 $q_j=e^{2\pi i(x_j+iy_j)/\epsilon}$；原文随后指出 $\epsilon\to0$ 对应大复结构极限，这是后文讨论 SYZ 猜想极限形式的伏笔。]

我们把这些环面纤维化记作 $f:X(B)\to B$ 与 $\check f:\check X(B)\to B$；它们显然是互相对偶的。

现在假设此外 $B$ 上还有一个 Hessian 度量 $g$，其局部势函数为 $K$。那么事实上 $X(B)$ 与 $\check X(B)$ 都成为 Kähler 流形：

\begin{proposition}{3.3.2}
$K\circ f$ 是 $X(B)$ 上的（局部）Kähler 势，它定义了 Kähler 形式 $\omega=2i\partial\bar\partial(K\circ f)$。该度量是 Ricci 平坦的，当且仅当 $K$ 满足实 Monge--Ampère 方程
\begin{equation*}
\det\frac{\partial^2K}{\partial y_i\partial y_j}=\text{常数}.
\end{equation*}
\end{proposition}

\begin{proof}
局部地取仿射坐标 $(y_i)$ 与复坐标 $z_j=\frac{1}{2\pi i}\log q_j=x_j+iy_j$，我们算得 $\omega=2i\partial\bar\partial(K\circ f)=\frac{i}{2}\sum\frac{\partial^2K}{\partial y_j\partial y_k}dz_j\wedge d\bar z_k$，它显然是正的。此外，若 $\Omega=dz_1\wedge\cdots\wedge dz_n$，则 $\omega^n$ 与 $\Omega\wedge\bar\Omega$ 成比例，当且仅当 $\det(\partial^2K/\partial y_j\partial y_k)$ 为常数。\proofend
\end{proof}

% [Tutor Note: 命题中写 $\omega=2i\partial\bar\partial(K\circ f)$，而证明中同一形式写作 $\frac{i}{2}\sum g^{jk}dz_j\wedge d\bar z_k$（后者见 p.10）；二者相差的是 $\partial\bar\partial$ 的归一化约定（$z_j=x_j+iy_j$ 下 $\partial\bar\partial$ 的系数），译文两处均照原文，未作统一。]

我们把该 Kähler 流形记作 $X(B,K)$。对偶地有

\begin{proposition}{3.3.3}
在 $T_B^*$ 的局部典范坐标 $y_i,\check x_i$ 下，$T_B^*$ 上的函数 $z_j=\check x_j+i\partial K/\partial y_j$ 在 $\check X(B)$ 上诱导出一个良定义的复结构，
% [原文 p.10]
使得典范辛形式 $\omega$ 是某个度量的 Kähler 形式。此外，该度量是 Ricci 平坦的，当且仅当 $K$ 满足实 Monge--Ampère 方程
\begin{equation*}
\det\frac{\partial^2K}{\partial y_j\partial y_k}=\text{常数}.
\end{equation*}
\end{proposition}

\begin{proof}
容易看出，坐标 $y_j$ 的仿射线性变换（以及坐标 $\check x_j$ 的相应适当变换）导致坐标 $z_j$ 的线性变换，因此它们诱导出一个在 $\check x_j\mapsto\check x_j+1$ 之下不变的、良定义的复结构，从而在 $\check X(B)$ 上诱导出复结构。然后计算得
\begin{equation*}
\omega=\sum d\check x_j\wedge dy_j=\frac{i}{2}\sum g^{jk}dz_j\wedge d\bar z_k
\end{equation*}
其中 $g_{ij}=\partial^2K/\partial y_j\partial y_k$。于是该度量是 Ricci 平坦的，当且仅当 $\det(g^{jk})=\text{常数}$，也当且仅当 $\det(g_{jk})=\text{常数}$。\proofend
\end{proof}

% [Tutor Note: 这里先出现 $\det(g^{jk})$ 再出现 $\det(g_{jk})$ 并非重复：$g^{jk}$ 是度量矩阵 $g_{jk}$ 的逆矩阵，二者行列式互为倒数，因此「等于常数」这一条件等价；文字层的 pdftotext 会丢掉上标而看起来像重复，须对照页图（p.10，700 dpi 裁图）才能正确抄写。]

和以前一样，我们把该 Kähler 流形记作 $\check X(B,K)$。这促使我们作如下定义。

\begin{definition}{3.3.4}
若局部势函数 $K$ 满足 Monge--Ampère 方程 $\det(\partial^2K/\partial y_i\partial y_j)=\text{常数}$，则称一个带有 Hessian 形式度量的仿射流形为 \emph{Monge--Ampère 流形}。
\end{definition}

Monge--Ampère 流形最早由 Cheng 与 Yau 在 [9] 中研究。

\begin{exercise}{3.3.5}
证明：由 Hessian 度量给出的 $T_B$ 与 $T_B^*$ 之间的等同，诱导出 Kähler 流形的典范同构 $X(B,K)\cong\check X(\check B,\check K)$，其中 $(\check B,\check K)$ 是 $(B,K)$ 的 Legendre 变换。
\end{exercise}

最后我们指出，可以在这一图景中引入 $B$-场。为使事情相对简单（从而不必过渡到广义复结构 [31]、[27]、[5]），我们把 $B$-场看作元素 $\mathbf{B}\in H^1(B,\Lambda_\R/\Lambda)$，其中 $\Lambda_\R=\Lambda\otimes_\Z\R$。注意到 $\Lambda_\R/\Lambda$ 在开集 $U$ 上的截面可以看作 $T_U/\Lambda|_U$ 的截面，这样的截面通过平移作用在 $T_U/\Lambda|_U$ 上，而这一作用关于标准的半平坦复结构实际上是全纯的。于是，代表 $\mathbf{B}$ 的 Čech 1-上链 $(U_{ij},\beta_{ij})$ 使我们能够通过在相交 $U_{ij}$ 上的平移重新粘合 $X(B)$。这给出一个新的复流形 $X(B,\mathbf{B})$。若此外还有一个定义度量的多值势函数 $K$，则这些平移保持度量，从而给出一个 Kähler 流形 $X(B,\mathbf{B},K)$。

于是镜像对称的完整玩具版本如下。数据由仿射流形 $B$ 连同势 $K$，以及 $B$-场 $\mathbf{B}\in H^1(B,\Lambda_\R/\Lambda)$、$\check{\mathbf{B}}\in H^1(B,\check\Lambda_\R/\check\Lambda)$ 组成。现在不难看出（如果你做过习题 3.3.5，你其实已经见过这一点）：用 $B$ 上仿射结构所定义出的格局部系统 $\check\Lambda$，与用 $\check B$ 上仿射结构所定义出的格局部系统 $\Lambda$，是同一个。于是我们说，对
\begin{equation*}
(X(B,\mathbf{B},K),\check{\mathbf{B}})
\end{equation*}
镜像于
\begin{equation*}
(X(\check B,\check{\mathbf{B}},\check K),\mathbf{B}).
\end{equation*}

% [Tutor Note: 原文此处把 affine structure 误拼作 "affine stucture"（缺 r；1200 dpi 裁图复核，p.10 对应行），不影响数学内容。]

这为镜像对称在一个简单情境中提供了相当令人满意的图景。镜像对称的更多方面可以在这一半平坦情境中展开，见 [40]。然而，归根结底，这只为一般情形提供了有限的洞见。以这种方式出现的、具有半平坦 Ricci 平坦度量的紧 Calabi--Yau 流形只能是复环面（Cheng 与 Yau 在 [9] 中证明）。

% [Tutor Note: 本块（b03）按「开始者写完」规则写到 p.10/p.11 跨界句完：即上面「……紧 Calabi--Yau 流形只能是复环面」一句整体属于 p.10 起头的句子，由本块译完；下一句「To deal with more interesting cases…」起于 p.11，归 tr02d。]
